\documentclass[10pt,a4paper]{amsart}
\usepackage[margin=19mm]{geometry}
\usepackage{amsmath,amssymb,mathtools}
\usepackage[scr=rsfs,cal=euler]{mathalfa}
\usepackage{xfrac}
\usepackage[unicode,hidelinks,bookmarks=false]{hyperref}
\newcommand{\ie}{i.e.\ }
\newcommand{\cf}{cf.\ }
\newcommand{\resp}{respectively\ }
\newcommand{\Subsection}{\subsection}
\providecommand{\nobreakdash}{\nobreak\hbox{-}}
\setlength{\emergencystretch}{2em}
\multlinegap0pt
\usepackage{bm}
\usepackage{bbm}
\usepackage{dsfont}
\usepackage{xspace}
\usepackage{cancel}
\usepackage{mathtools}
\usepackage{amsthm}
\makeatletter
\@ifundefined{proposition}{\newtheorem{proposition}{Proposition}}{}
\makeatother
\title[Beyond DSA: Conjugacy-based Comparison of Dynamical Systems]{Beyond DSA: Conjugacy-based Comparison of Dynamical Systems}
\author{Prakhar Godara, Pang Shiang Tay, Marcelo G. Mattar}
\date{Source: arXiv:2607.04493v1 (CC BY 4.0)}
\begin{document}
\maketitle

\noindent\emph{Source and licence.} Beyond DSA: Conjugacy-based Comparison of Dynamical Systems. Version arXiv:2607.04493v1 (CC BY 4.0), distributed under the Creative Commons Attribution 4.0 International licence, as recorded in the arXiv licence metadata for this exact version and shown on the abstract page https://arxiv.org/abs/2607.04493v1. Full rights notice: licenses/arxiv-2607.04493-CC-BY-4.0.txt.\par\smallskip

\noindent\textit{Editorial scope.} Counted unit: the Proposition whose source label is \texttt{prop:bernoulli\_counterexample\_appendix} in the article (source0.tex lines 473--495, source label \texttt{prop:bernoulli\_counterexample\_appendix}), delivered together with the article's own complete original proof of it (source0.tex lines 497--556, one \texttt{proof} environment of 60 lines). No mathematics is written, repaired or re-derived: statement and proof are the article's own text line for line. Required context retained uncounted: none. Every definition, display and label of those lines is retained unchanged. Omitted from this package and not redistributed: 1--472, 496--496, 557--1480. Nothing inside a retained excerpt is omitted or altered: every retained body is delivered exactly as the article's lines are written, so no omission receipt of any kind accompanies this package. Citations: the retained excerpts carry no citation command at all, so no bibliography entry of the article is transcribed and no reference is redistributed. Reference adaptation: none was needed and none was made; every reference inside the delivered unit resolves inside it, so no unresolved reference or citation is produced. Printed slips kept literal: no printed word, symbol, hypothesis or number of the counted statement or of its proof was altered. Layout and class: the article's own class is \texttt{article} with packages including neurips\_2026, enumitem, natbib, graphicx, wrapfig, subfig, subcaption, dcolumn, inputenc, blindtext, xcolor, algorithm, algpseudocode, tikz; the wrapper is the lane's pinned \texttt{amsart} editorial wrapper at 10pt on A4, which restates the theorem-like environments the retained text needs and adds the article's own macro definitions that the retained text needs (none), with extra wrapper packages (bm, bbm, dsfont, xspace, cancel, mathtools). The delivered numbering is the unit's own. Assets: no retained span contains a figure environment, a graphics-inclusion command, a PDF-page inclusion, a table or a TikZ picture; no image file is copied into this package, the package declares no asset dependency and no image file is redistributed with it. Licence: version arXiv:2607.04493v1 of this article is distributed under the Creative Commons Attribution 4.0 International licence, as recorded in the arXiv licence metadata for this exact version and shown on the abstract page https://arxiv.org/abs/2607.04493v1. A full rights notice is delivered at licenses/arxiv-2607.04493-CC-BY-4.0.txt. Characters: every retained excerpt is pure ASCII, so the delivered engine, XeLaTeX with its default Latin Modern text font, reports no missing character.
\begin{proposition}
\label{prop:bernoulli_counterexample_appendix}
Let
\[
X:=\{0,1\}^{\mathbb Z}\times \{0,1\},
\qquad
Y:=\{0,1,2\}^{\mathbb Z}\times \{0,1\},
\]
equipped with the product probability measures
\[
\mu:=\mu_2\times \eta,
\qquad
\nu:=\mu_3\times \eta,
\]
where \(\mu_2\) and \(\mu_3\) are the Bernoulli product measures on
\(\{0,1\}^{\mathbb Z}\) and \(\{0,1,2\}^{\mathbb Z}\), respectively, and \(\eta\) is the uniform measure on \(\{0,1\}\). Define
\[
T:=\sigma_2\times \mathrm{id},
\qquad
S:=\sigma_3\times \mathrm{id},
\]
where \(\sigma_2\) and \(\sigma_3\) are the two-sided Bernoulli shifts. Then the Koopman operators \(\mathcal K_T\) and \(\mathcal K_S\) are unitarily equivalent, but \(T\) and \(S\) are not conjugate.
\end{proposition}

\begin{proof}
We prove the two claims separately.

First, \(\mathcal K_T\) and \(\mathcal K_S\) are unitarily equivalent. For \(k=2,3\), let
\[
X_k:=\{0,1,\dots,k-1\}^{\mathbb Z},
\qquad
(\sigma_kx)_n=x_{n+1},
\]
with \(\mu_k\) the corresponding Bernoulli product measure. It is a standard fact that
\[
L^2(X_k,\mu_k)
=
\mathbb C \mathbf 1
\oplus
\bigoplus_{j=1}^\infty \mathcal H_j^{(k)},
\]
where on each \(\mathcal H_j^{(k)}\), the Koopman operator \(\mathcal K_{\sigma_k}\) acts as a copy of the bilateral shift on \(\ell^2(\mathbb Z)\). Since this decomposition is the same for \(k=2\) and \(k=3\), there exists a unitary
\[
W:L^2(X_2,\mu_2)\to L^2(X_3,\mu_3)
\]
such that
\[
W\mathcal K_{\sigma_2}=\mathcal K_{\sigma_3}W.
\]

Now
\[
L^2(X,\mu)\cong L^2(X_2,\mu_2)\oplus L^2(X_2,\mu_2),
\qquad
L^2(Y,\nu)\cong L^2(X_3,\mu_3)\oplus L^2(X_3,\mu_3),
\]
and under these identifications,
\[
\mathcal K_T=\mathcal K_{\sigma_2}\oplus \mathcal K_{\sigma_2},
\qquad
\mathcal K_S=\mathcal K_{\sigma_3}\oplus \mathcal K_{\sigma_3}.
\]
Hence \(\widetilde W:=W\oplus W\) is unitary and satisfies
\[
\widetilde W \mathcal K_T = \mathcal K_S \widetilde W.
\]
Thus \(\mathcal K_T\) and \(\mathcal K_S\) are unitarily equivalent.

Next, \(T\) and \(S\) are not conjugate. If there were a bijection \(h:X\to Y\) with
\[
h\circ T=S\circ h,
\]
then \(T\) and \(S\) would have the same number of periodic points of each period. But
\[
\#\mathrm{Fix}(\sigma_k^n)=k^n,
\]
since an \(n\)-periodic bi-infinite sequence is determined by a word of length \(n\). Therefore
\[
\#\mathrm{Fix}(T^n)=2\cdot 2^n,
\qquad
\#\mathrm{Fix}(S^n)=2\cdot 3^n.
\]
These are unequal for every \(n\ge 1\), so no such \(h\) can exist. Hence \(T\) and \(S\) are not conjugate.
\end{proof}

\end{document}
